\section{Conclusion}\label{sec_conclusion}

We introduced in this paper new time-scale-frequency transforms that can adapt their time-frequency resolution to the analyzed signal, through the frequency domain and time domain focus functions $f\to\sigma^\nu_f$ and $f\to\sigma^\tau_f$. Based upon short time Fourier transform or continuous wavelet transform, the proposed transforms adapt dynamically the scale/bandwidth of analysis windows or wavelet as a function of frequency or time, leading to non-linear transforms. Under suitable assumptions on focus functions, we could prove first important results on the transforms such as the well-definedness on $L^2$, and norm controls similar to the one obtained in the linear case.

In Theorems~\ref{theo_bounding_time_marseillan_transform} and~\ref{theorem_norm_control_frequency_focus}, we obtain a control of the type
\[
\sup_{\substack{f\in L^2(\R)\\ f\ne 0}} \frac{\|Mf\|_{L^2}}{\|f\|_{L^2}} = \sup_{\substack{f\in L^2(\R)\\ f\ne 0}} C_f^{1/2}\ ,
\]
where $C_f$ depends on $f\in L^2(\R)$ only through the focus function $\sigma_f$. More specific assumptions on the focus functions are needed to insure a finite upper bound for $C_f$ with $f\in L^2$. For example, one may specify that $1\le\sigma_f\le\sigma_\text{max}$ for some prescribed $\sigma_\text{max}$, as we did in numerical illustrations. It would be  interesting to study more thoroughly generic mappings $f\to\sigma_f$ and derive sufficient conditions insuring the finiteness of $C_f$.

Note that, due to the non-linearity of the transform, the above quantity doesn't define a norm for the transform. Lipschitz continuity, \ie\, the existence of a constant $C(M)$ such that
\begin{equation}\label{eq_lipchitz_constant}
\forall f_1,f_2\in L^2\ ,\quad \|Mf_1 - Mf_2 \|_{L^2} \leqslant C(M)\|f_1-f_2\|_{L^2}
\end{equation}
would clearly be of interest too (see also generalized operator norms introduced and studied in~\cite{wei2020development}).
Proving the existence of such Lipschitz constants for $M^\nu$ and  $M^\tau$ would ensure their uniform continuity. 
%We saw in Lemma~\ref{lemma_density_result} that the time-focus transform has a local continuity property (under some hypothesis) and  Proposition~\ref{prop_continuity_kernel} can in a certain way give us the local continuity of the frequency-focus transform.
However, none of the two above mentioned results is strong enough to prove the existence of a Lipschitz constant that satisfies Equation~\eqref{eq_lipchitz_constant}.
We plan to follow this line in the near future.

Of interest too for the inversion of the non-linear transforms would be to investigate which conditions would guarantee the existence of a a constant $c(M)$ such that
\begin{equation}
    \forall f_1,f_2\in L^2(\R)\ ,\quad  c(M)\|f_1-f_2\|_{L^2} \leqslant\|Mf_1 - Mf_2 \|_{L^2}\ .
\end{equation}
Such property would guarantee injectivity of the non-linear transform. Again, the lower bounds provided in Theorems~\ref{theo_bounding_time_marseillan_transform} and~\ref{theorem_norm_control_frequency_focus} are not sufficient to yield directly injectivity, even though the bound does not depend on the analyzed function $f$.

\begin{comment}
cannot be understood in terms of operator norms. Exploring the ways to define norms on general map on Banach spaces (in our case $L^2$) would be an interesting perspective. A first way could be to follow~\cite{wei2020development}, where generalized operator norms are introduced, which cover the non-linear case. For $F:L^2(\R)\rightarrow L^2(\R^2)$, define
\begin{equation}
    \rho(F) := \max\left\{   \sup_{f\not=0}\frac{\|F(f)\|_{L^2}}{\|f\|_{L^2}}, \|F(0)\|_{L^2}  \right\} \in [0,+\infty]\ .
\end{equation}
%Such way to define a norm insist on the bounded aspect of the application. 
It is quite easy to see from theorem~\ref{theo_bounding_time_marseillan_transform} that the time-focus transform $M^\tau$ satisfies
\begin{equation}
    \rho(M^\tau) \leqslant \|h\|_{L^\infty}|\supp(h)|^{1/2}\ ,
\end{equation}
and has therefore a finite norm in the above sense. However, at this point, the control from Theorem~\ref{theorem_norm_control_frequency_focus} does not allow us to conclude a similar control for $\rho(M^\nu)$.

An alternative is to define the norm of $F$ as its Lipschitz constant (since continuity and boundedness coincide in the case of linear operators), \ie\, the existence of a constant $C(F)$ such that
\begin{equation}\label{eq_lipchitz_constant}
\forall f_1,f_2\in L^2\ ,\quad \|F(f_1) - F(f_2) \|_{L^2(\R^2)} \leqslant C(F)\|f_1-f_2\|_{L^2(\R)}\ .
\end{equation}
Proving the existence of such Lipschitz constants for $M^\nu$ and  $M^\tau$ would ensure their uniform continuity. We saw in Lemma~\ref{lemma_density_result} that the time-focus transform has a local continuity property (under some hypothesis) and  Proposition~\ref{prop_continuity_kernel} can in a certain way give us the local continuity of the frequency-focus transform. However, none of the two results is strong enough to prove the existence of a Lipschitz constant that satisfies Equation~\eqref{eq_lipchitz_constant}.
We plan to study this question in the future.

We also plan to investigate in which conditions there is a constant $c(M)$ such as
\begin{equation}
    \forall f_1,f_2\in L^2\ ,\quad  c(M)\|f_1-f_2\|_{L^2(\R)} \leqslant\|Mf_1 - Mf_2 \|_{L^2(\R^2)}\ .
\end{equation}
Such property would guarantee injectivity of the non-linear transform.
\end{comment}


\smallskip
A main further goal will be to study the invertibility of such non-linear transforms. From our results, inverse transforms can be obtained if both the transform $Mf$ and the focus function $\sigma_f$ are known, but not in situations where only the transform $Mf$ is known. A first step would be to analyze in which conditions an approximate inverse can be obtained when an approximation of the focus function is available. The above-mentioned problems are likely to play a role for this question. This may open the door to iterative inversion methods.

\smallskip
Last but not least, we plan to head to concrete applications of this approach, in particular in the context of audio perception modelling, which was one of the main motivations for this work. For that, we plan to investigate further focus functions that could be relevant in applications, starting from the simple models and examples described in Section~\ref{se:numerical.experiments}, and study more thorough applications to real signals.
%Those results lead us to further purposes. The  first one would be to study the definitions of the focus function and to understand which properties of the signals we want to focus on. Another goal would be to estimate those focus function : for a given signal how to estimate with a certain precision $\eps$ the focus function associated to $f$. Having such estimate focus function we can study a linearized transform : instead of using the proper focus function of the signal, use the estimate focus function. Hence the obtained transform is linear and try to control the error with the main operator. Then consequence of the theorem \ref{theorem_norm_control_frequency_focus} and \ref{theo_bounding_time_marseillan_transform} is that for fixed focus functions ($\sigma_f^\nu$ and $\sigma_f^\tau$) the linearized transforms are injective which lead us to the idea of inversion of the linearized operator.